\section{A conjecture: rational-circle rigidity}\label{sec:rigidity}

Everything in this paper concerns the \emph{rational-circle relaxation}: an integer frequency matrix $A\in\Z^{4\times r}$
satisfying the \emph{full-depth condition} that in every column at least three rows attain the column's maximal
absolute value (the condition Lemma~\ref{lem:depth} forces for every primitive square; for a squarefree center it says
that every coordinate is nonzero in at least three rows of a matrix with entries in $\{0,\pm1\}$), together with the
question of whether a point of $T(\Q)^r$ is admissible for $A$ in the sense of Section~\ref{sec:introduction}. Two equations in $r$ parameters: the solution set has
dimension $r-2$, finite for two parameters, curves for three, and surfaces for $r\ge4$. The results above suggest a
single statement.

\begin{conjecture}[Rational-circle rigidity]\label{conj:rigidity}
For every $r\ge2$ and every full-depth configuration $A$, there is no admissible point of $T(\Q)^r$ for $A$.
\end{conjecture}

The conjecture is strictly stronger than the nonexistence of a primitive magic square of nine distinct squares: every
such square yields a full-depth configuration and an admissible point of $T(\Q)^r$ (Lemma~\ref{lem:depth} and~\eqref{eq:areaangle}),
and the conjecture discards the Gaussian-prime condition entirely. Theorem~\ref{thm:main-two} proves it for two parameters
through exponent sum seven and Corollary~\ref{cor:parallelogram} for every parallelogram configuration, over the reals; for two parameters at all exponents the conjecture would give a purely trigonometric proof of the archive's theorem
that no primitive square has a two-prime central root (Section~\ref{sec:prior}), a theorem presently proved only through the
Gaussian-prime arithmetic;
Conjecture~\ref{conj:points} is, by Theorem~\ref{thm:main-d1}, equivalent to its instance for the configuration D1, and by
Corollary~\ref{cor:fourtargets} implies its instances for B2, C and E2. If the conjecture fails, the smallest $r$ at which it
fails locates where the arithmetic of the circle points becomes essential: below that $r$ the obstruction to magic squares
is trigonometric, and from that $r$ on it is arithmetic. The three-parameter cases being open, the first unexplored test is
$r=4$, where the solution sets are surfaces. The two results of this section concern that case.

\subsection{Four parameters at bounded height}

Write $\tan(\chi_j/2)=[p_j:q_j]\in\PP(\Q)$ in lowest terms, with height $\max(|p_j|,|q_j|)$; thus $0=[0:1]$ and
$\infty=[1:0]$ have height one, and correspond to $\chi_j=0$ and $\chi_j=\pi$.

\begin{proposition}\label{prop:four}
The $11{,}862$ four-element sets of nonzero vectors of $\{0,\pm1\}^4$ modulo sign that satisfy the full-depth condition
form $113$ orbits under signed coordinate permutations. No such configuration admits an admissible point of $T(\Q)^4$ at
which three of the four parameters $\tan(\chi_j/2)$ have projective height at most $4$.
\end{proposition}

\begin{proof}[Computer-assisted proof]
For each orbit representative, each choice of a coordinate $j$, and each triple of values of the other parameters among the
$24$ points of $\PP(\Q)$ of height at most $4$ (including $0$ and $\infty$), every row's sine is an integer linear form in
$(\cos\chi_j,\sin\chi_j)$ once the positive common denominator is cleared. Each of the $48$ signed systems of
Lemma~\ref{lem:roles} is two such linear equations, and three cases arise. If their determinant is nonzero, the unique
solution is tested exactly against the unit circle. If the determinant vanishes but the equations are not both identically
zero, the consistent cases reduce to a line meeting the circle, decided by an integer square-root test. If both equations
vanish identically — which happens at $12{,}240$ of the $113\cdot4\cdot13{,}824\cdot48$ system evaluations, at parameter
triples containing torsion values — the relations hold for \emph{every} value of $\chi_j$, and the four sines
$A_rc+B_rs+C_r$ form a one-parameter family; such a family contains admissible points (all but finitely many) if and only if
no coefficient triple $(A_r,B_r,C_r)$ vanishes and no two are equal up to sign, which is tested exactly, and in every
occurrence the family is degenerate. Surviving points are tested for admissibility exactly. Since every admissible point has
some three parameters in the grid for some ordering of the coordinates, and every configuration is a signed permutation of a
representative, the sweep is complete for the stated height. It finds no admissible point; height $3$ was swept separately
with the same outcome, and an earlier sweep of the finite chart at height $5$ (without $0,\infty$ and without the family
branch) also found none. The code (\path{verification/four_prime/r4search.py}), its checkpoints, and unit tests —
reconstruction of all four sines from the linear forms at $400$ random rational points, recovery of planted circle points,
and the admissibility and singular branches — accompany the paper.
\end{proof}

A heuristic count predicts the opposite: a surface in four rational parameters should carry many rational points. The
next result explains the absence for the most symmetric configuration.

\subsection{The Hadamard solution surfaces are of general type}

Let $A_{\mathrm H}$ be the Hadamard configuration with rows $(1,1,1,1),(1,1,-1,-1),(1,-1,1,-1),(1,-1,-1,1)$, and put
\[
\alpha=\chi_1+\chi_2,\quad \beta=\chi_3+\chi_4,\quad \gamma=\chi_1-\chi_2,\quad \delta=\chi_3-\chi_4,\qquad
u=\gamma+\delta,\quad v=\gamma-\delta,
\]
so that $(s_0,s_1,s_2,s_3)=(\sin(\alpha+\beta),\sin(\alpha-\beta),\sin(\gamma+\delta),\sin(\gamma-\delta))$. In every signed
system~\eqref{eq:roles} the two equations can be solved for $(s_2,s_3)$ as linear forms in $(s_0,s_1)$ — the coefficient
block on $(s_2,s_3)$ is invertible in all three index patterns — so each system reads
\begin{equation}\label{eq:bidouble}
\sin u=F_1(\alpha,\beta),\qquad \sin v=F_2(\alpha,\beta),
\end{equation}
with $F_1,F_2\in\Q\sin(\alpha+\beta)+\Q\sin(\alpha-\beta)$. Up to the symmetries $u\mapsto-u$, $v\mapsto-v$,
$u\leftrightarrow v$ and $\alpha\leftrightarrow\beta$, exactly four pairs $(F_1,F_2)$ occur:
\[
\begin{array}{cll}
\text{(I)} & \sin\alpha\cos\beta, & \cos\alpha\sin\beta;\\
\text{(II)} & 2\cos\alpha\sin\beta, & 2\sin\alpha\cos\beta;\\
\text{(III)} & 2\cos\alpha\sin\beta, & 3\cos\alpha\sin\beta-\sin\alpha\cos\beta;\\
\text{(IV)} & 2\cos\alpha\sin\beta, & \sin\alpha\cos\beta+3\cos\alpha\sin\beta .
\end{array}
\]

\begin{theorem}[All Hadamard solution surfaces have general type]\label{thm:gentype}
For each of the $48$ signed systems of $A_{\mathrm H}$, let $S\subset(\PP)^4$ be the closure of its solution set in the
coordinates $\tan(\chi_j/2)$, and let $S''\subset(\PP)^4$ be the closure of~\eqref{eq:bidouble} in the coordinates
$\tan(\alpha/2),\tan(\beta/2),\tan(u/2),\tan(v/2)$. Then:
\begin{enumerate}
\item $S''$ is the irreducible $(\Z/2)^2$-cover of $\PP\times\PP$ branched along $B_1=\{F_1=\pm1\}$ and
$B_2=\{F_2=\pm1\}$, each of bidegree $(4,4)$, reduced, with no common component. Every component of $B_1\cup B_2$ has only
ordinary nodes as singularities, the two components of each $B_i$ meet transversally, no three components pass through
one point, the nodes of $B_1$ lie off $B_2$ and conversely, and every point of $B_1\cap B_2$ is either transversal or a
simple tangency.
\item Consequently $S''$ has only rational double points of type $A_1$, and its minimal resolution is a minimal surface of
general type with
\[
p_g=11,\qquad q=0,\qquad K^2=32,\qquad \chi(\mathcal O)=12 .
\]
\item The change of angles is an isogeny of tori of degree $8$ and induces a dominant, generically finite rational map
$S\dashrightarrow S''$ of degree $8$ (the induced map of the toric compactifications need not be a morphism), under which
every two-dimensional component of $S$ dominates $S''$; hence every such component has Kodaira dimension~$2$.
\end{enumerate}
\end{theorem}

\begin{proof}
(1) In half-angle coordinates $(a,b,w,y)$, \eqref{eq:bidouble} becomes two equations of multidegrees $(2,2,2,0)$ and
$(2,2,0,2)$; the first is the double cover $F_1w^2-2w+F_1=0$ of the $(a,b)$-plane, branched where its discriminant
$4(1-F_1^2)$ vanishes, i.e.\ on $B_1$, and likewise for $B_2$. Each $\{F_i=\pm1\}$ clears to a polynomial of bidegree
$(2,2)$. The fibre product of two irreducible double covers with distinct branch loci is irreducible. The geometry of the
branch curves is a finite Gr\"obner-basis computation in all four affine charts of $\PP\times\PP$
(\path{verification/four_prime/branchcheck.py}, \path{branchfine.py}), with the following outcome. Type~(I): each $(2,2)$
component splits over $\Q(i)$ into two smooth $(1,1)$ curves, e.g.\ $(1+a^2)(1+b^2)-2a(1-b^2)=(1-a)^2+b^2(1+a)^2$, so
$B_i$ is four rational curves with $2+2+8=12$ transversal nodes, and $B_1\cap B_2$ consists of $32$ transversal points.
Type~(II): all four components are smooth curves of genus one; $B_1$ and $B_2$ each have $8$ transversal nodes, and
$B_1\cap B_2$ consists of $16$ simple tangencies and no transversal points — the full intersection number
$(4,4)\cdot(4,4)=32$ — at $\tan(\alpha/2),\tan(\beta/2)\in\{\pm1\pm\sqrt2\}$, i.e.\ at $\alpha,\beta\in\{\pm\pi/4,\pm3\pi/4\}$,
where $\cos2\alpha=\cos2\beta=0$. Types~(III), (IV): the components of $B_1$
are smooth of genus one, the components of $B_2$ split over $\Q(\sqrt2)$ into smooth $(1,1)$ curves (the coefficients
$3\pm2\sqrt2=(1\pm\sqrt2)^2$ appear), all intersections within each $B_i$ are transversal nodes, the nodes lie off the other
divisor, there are no triple points, and $B_1\cap B_2$ consists of $16$ transversal points and $8$ simple tangencies
($16+2\cdot8=32$). In each type the counts were confirmed against the intersection number.

(2) Local analysis of the cover $z^2=f$, $w^2=g$ over a point of the base: over a smooth point of one branch curve not on
the other it is smooth; over a node of $B_1$ (resp.\ $B_2$) it is $z^2=xy$ on each of two sheets, an $A_1$ point; over a
transversal point of $B_1\cap B_2$ it is $z^2=x$, $w^2=y$, smooth; and over a simple tangency, with $f=y$, $g=y-x^2$, it is
$w^2-z^2+x^2=0$, again an $A_1$ point. Hence all singularities of $S''$ are canonical. The cover is a normal bidouble cover
with branch data $D_1=B_1$, $D_2=B_2$, $D_3=0$; by Pardini's structure theorem for normal abelian covers~\cite{Pardini},
$\pi_*\mathcal O_{S''}=\bigoplus_i\mathcal O(-L_i)$ with $2L_1=D_2$, $2L_2=D_1$, $2L_3=D_1+D_2$, so $L_1=L_2=(2,2)$ and
$L_3=(4,4)$. Rational double points are rational and crepant, so the minimal resolution has
\[
p_g=\sum_{i=1}^{3}h^0\bigl(K_X+L_i\bigr)=h^0(\mathcal O)+h^0(\mathcal O)+h^0(\mathcal O(2,2))=1+1+9=11,\qquad
q=\sum_{i=1}^{3}h^1(K_X+L_i)=0,
\]
\[
K^2=4\bigl(K_X+\tfrac12(D_1+D_2)\bigr)^2=4\,(2,2)^2=32 .
\]
Independently, adjunction on the complete intersection gives $K_{S''}=\mathcal O(2,2,0,0)|_{S''}$, the pullback of an ample class
under a finite map and hence ample, with $K^2=(2h_1+2h_2)^2(2h_1+2h_2+2h_3)(2h_1+2h_2+2h_4)=32$, and the Koszul resolution
gives $\chi(\mathcal O_{S''})=1-(-1)-(-1)+9=12=1-q+p_g$. Ampleness of $K$ on the singular model shows the resolution is
minimal; Noether's inequality $2p_g-4\le K^2\le 9\chi$ reads $18\le32\le108$~\cite{BHPV}.

(3) The change of angles $\chi\mapsto(\alpha,\beta,\gamma,\delta)$ has determinant $4$ and $(\gamma,\delta)\mapsto(u,v)$
has determinant $2$; on the tori this is an isogeny of degree $8$, so the induced rational map $S\dashrightarrow S''$ is
dominant and generically finite of degree $8$, and every surface component of $S$ dominates the irreducible surface $S''$.
Kodaira dimension does not decrease under a dominant generically finite rational map (resolve the indeterminacy and apply
the statement for morphisms).
\end{proof}

\begin{corollary}
Under the Bombieri--Lang conjecture~\cite{Lang}, every component of every Hadamard solution surface has only finitely many
rational points outside finitely many curves. This is consistent with Proposition~\ref{prop:four} and with
Conjecture~\ref{conj:rigidity}; it proves neither.
\end{corollary}

\subsection{Two further hypotheses}

Conjecture~\ref{conj:rigidity} asserts emptiness, which is far stronger than what the geometry suggests: for a surface or
higher-dimensional variety of general type, the Bombieri--Lang conjecture predicts only that rational points are not
Zariski dense. We therefore record the statement that the geometry supports directly, and a structural hypothesis about
where the rational points that do exist are located.

\begin{conjecture}[Weak form]\label{conj:weak}
For every full-depth configuration, the admissible points of $T(\Q)^r$ are not Zariski dense in any component of the
solution variety. By Bombieri--Lang this holds for every component of general type, in particular for every component in
Theorem~\ref{thm:gentype}.
\end{conjecture}

Every rational point of a solution variety found in this paper lies on a \emph{torsion coset}: the eight points of $K$
on $\mathcal C$ have $\alpha,\beta,\gamma\in\{0,\pm\pi/2,\pi\}$; the trivial sieve classes in the extended record are
torsion points; and all $12{,}240$ free-angle families in the height-$4$ sweep occur at parameter triples containing a
torsion value. The torsion of $T(\Q)$ is $\{\pm1,\pm i\}$, and a rational point of $T(\Q)^r$ whose coordinates are
multiplicatively independent lies on no proper subtorus coset, so the following is a genuine restriction.

\begin{conjecture}[Torsion-coset hypothesis]\label{conj:torsion}
For every full-depth configuration, every rational-circle point of its solution variety satisfies a nontrivial relation
$\prod_j z_j^{\,n_j}\in\{\pm1,\pm i\}$ with $(n_j)\in\Z^r\setminus\{0\}$; that is, the rational points lie on
finitely many torsion cosets of proper subtori.
\end{conjecture}

The hypothesis has an inductive consequence. If the relation has some coefficient $n_j=\pm1$, eliminating $\chi_j$
turns the configuration into one in $r-1$ parameters, with integer entries of larger absolute value in general, so
Conjecture~\ref{conj:rigidity} for $r$ would follow from Conjecture~\ref{conj:torsion} for $r$ together with
Conjecture~\ref{conj:rigidity} for $r-1$; relations with all $|n_j|\ge2$ instead pass to a finite cover. In this sense
the torsion-coset hypothesis would reduce rigidity in all dimensions to the three-parameter cases, of which D1 — the
completeness of $K$ — is the first. It is also the mechanism by which a search finds nothing: the rational points that
the equations do admit sit on degenerate loci where two sine magnitudes coincide identically. We offer it as the
hypothesis whose confirmation or refutation would most change the picture.

\subsection{Evidence for the torsion-coset hypothesis and a degeneracy-rank lemma (22 October 2026)}\label{sec:torsionprobe}
The sweep of Proposition~\ref{prop:four} was re-instrumented to record \emph{every} exact solution of the $48$ signed systems, admissible or not, for all $113$ orbits, all four choices of solved coordinate, and all triples of the other parameters of projective height at most $H$, and each solution was tested for Conjecture~\ref{conj:torsion} by a complete criterion: with $\tan(\chi_j/2)=p_j/q_j$ one has $z_j=(q_j+ip_j)/(q_j-ip_j)$, whose image in $\Q(i)^\times/\{\pm1,\pm i\}$ is determined by the Gaussian-prime factorisation of $q_j+ip_j$; a relation $\prod z_j^{n_j}\in\{\pm1,\pm i\}$ exists if and only if the four exponent vectors are linearly dependent, and the rank of these vectors (the \emph{exponent rank}) is computed exactly.

\begin{proposition}\label{prop:torsiondata}
At $H=3$ there are $273{,}024$ distinct rational points of the solution varieties, at $H=4$ there are $662{,}528$, and at $H=5$ there are $1{,}916{,}672$ (each grid contains the previous one). Every one of them satisfies a nontrivial relation $\prod z_j^{n_j}\in\{\pm1,\pm i\}$; none is admissible; every one has a vanishing sine and two sines of equal magnitude. The exponent rank is at most $2$ at $H=3$; at $H=4$ it equals $3$ for $13{,}824$ points and at $H=5$ for $235{,}008$ points, in the same two orbits.
\end{proposition}
The rank-three points refute the sharper statement, suggested by the $H\le3$ data, that every rational point lies on a torsion coset of a subtorus of dimension at most two. They are explained by the following elementary lemma.

\begin{lemma}[Degeneracy relations]\label{lem:degrel}
Let $z\in T(\Q)^r$ and $s_j=\sin(A_j\cdot\chi)$. If $s_r=0$ then $\prod_j z_j^{A_{rj}}=\pm1$. If $|s_i|=|s_k|$ with $i\ne k$ then $\prod_j z_j^{A_{ij}-A_{kj}}=\pm1$ or $\prod_j z_j^{A_{ij}+A_{kj}}=\pm1$, with a nonzero exponent vector because the rows are distinct projective vectors. Consequently every rational point with a vanishing sine or two sines of equal magnitude satisfies Conjecture~\ref{conj:torsion}; a point with both degeneracies has exponent rank at most $3$, with equality only if $A_i\pm A_k$ is a rational multiple of the vanishing row $A_r$.
\end{lemma}
\begin{proof}
$s_r=0$ means $A_r\cdot\chi\in\pi\Z$, and $|s_i|=|s_k|$ means $A_i\cdot\chi\equiv\pm A_k\cdot\chi\pmod\pi$; exponentiate. Two independent relations lower the rank to at most $2$.
\end{proof}
\begin{theorem}[Trichotomy on a signed system]\label{thm:trichotomy}
Let $z\in T(\Q)^r$ satisfy a signed system~\eqref{eq:roles}, $s_i=\varepsilon s_k+\eta s_l$, $s_j=\varepsilon s_k-\eta s_l$, and put $A=\max(|s_k|,|s_l|)$, $B=\min(|s_k|,|s_l|)$. Then $\{|s_i|,|s_j|\}=\{A+B,A-B\}$, and exactly one of the following holds: (i) some sine vanishes; (ii) $z$ is admissible; (iii) $A=2B>0$. In case (i) two of the remaining sines have equal magnitude, so by Lemma~\ref{lem:degrel} the point lies on a torsion coset and has exponent rank at most $3$.
\end{theorem}
\begin{proof}
For real $x,y$ the magnitudes of $x+y$ and $x-y$ are $A+B$ and $A-B$ with $A=\max(|x|,|y|)$, $B=\min(|x|,|y|)$. If no sine vanishes then $B>0$ and $A\ne B$, and the four magnitudes $A,B,A+B,A-B$ are pairwise distinct unless $A-B=B$; by Lemma~\ref{lem:roles} the point is admissible exactly when they are distinct. For (i): if $s_l=0$ then $s_i=s_j$; if $s_k=0$ then $|s_i|=|s_j|=|s_l|$; if $s_i=0$ or $s_j=0$ then $|s_k|=|s_l|$.
\end{proof}
Case (iii) is the only way a rational point of a solution variety can avoid both admissibility and a monomial relation, and it is not a torsion-coset condition. Every one of the $1{,}916{,}672$ points recorded at $H\le5$ (each grid contains the previous one, so the sweeps are not added) is of type (i); a direct classification of every point found no point of type (iii). We therefore state the hypothesis that the data are consistent with and that implies both Conjecture~\ref{conj:rigidity} and Conjecture~\ref{conj:torsion}; it is strictly stronger than Conjecture~\ref{conj:rigidity}, since it also excludes type (iii), and the evidence for the extra exclusion is only the absence of type (iii) points at $H\le5$, where the grid carries the split primes $5$, $13$, $17$, $29$ and $41$:
\begin{conjecture}[Vanishing-sine rigidity]\label{conj:vanishing}
For every full-depth configuration, every rational point of every signed system has a vanishing sine.
\end{conjecture}
Conjecture~\ref{conj:vanishing} is equivalent to Conjecture~\ref{conj:rigidity} together with the nonexistence of rational points of type (iii). The next computation is therefore not a larger admissibility sweep but a direct search for type (iii) points: rational-circle solutions of a signed system with $|s_l|=2|s_k|$ or $|s_k|=2|s_l|$ and all sines nonzero.

The parallel-triple condition of the lemma holds for exactly $23$ of the $113$ orbit representatives, and the two orbits that carry rank-three points at $H=4$ and at $H=5$ are among them; every recorded rank-three point has a vanishing row parallel to the sum or difference of its two equal-magnitude rows. The remaining $21$ orbits are candidates for larger heights; the criterion is necessary, not sufficient. The lemma also shows that Conjecture~\ref{conj:rigidity} implies Conjecture~\ref{conj:torsion} for every solution-variety point that fails admissibility through a vanishing sine or coinciding magnitudes; the only other failure mode, $a=2b$, yields no monomial relation and was never observed. The code, the per-orbit counts and the twenty recorded rank-three examples are in the companion repository.

Theorem~\ref{thm:gentype} shows that the general-type behaviour Bremner~\cite{Bremner2001} and
Bruin--Thomas--V\'arilly-Alvarado~\cite{BTV} established for the magic-square surface is already present in the trigonometric
relaxation, with the Gaussian-prime arithmetic removed. The decoupling that makes the proof elementary — the rows pair into
$\{v_0,v_1\}$ and $\{v_2,v_3\}$ with $v_0\pm v_1$ and $v_2\pm v_3$ supported on complementary coordinate pairs — is special
to the Hadamard configuration. For a general full-depth configuration the two relations are divisors of multidegree $(2,2,2,2)$ in $(\PP)^4$; whenever
they have no common component the solution surface is their complete intersection, adjunction gives
$K_S=\mathcal O(2,2,2,2)|_S$, which is ample, and the surface is of general type whenever its singularities are canonical. Verifying that for the other $112$ orbits, and
locating any configuration whose surface is instead rational — which would be a direct route to a rational-circle solution
and would refute Conjecture~\ref{conj:rigidity} — is the natural continuation.
